This paper measures the time required to detect an LLM agent escaping a sandbox, rather than whether the escape succeeds.
Using Sentinel, a containerized harness, we evaluate 416 runs on four machines against three telemetry signals that never query the model: process liveness, persistent canary tokens, and outbound traffic.
Detection lands in 2.22 seconds, and nothing acts for a further 44.
Three findings follow.
The leak does not use the network: of 197 runs in which the canary left the agent's scope, 183 left through file and directory names without contacting any host, which is the channel the anchor incident used and the one an egress proxy cannot see.
The delivery channel dictates compliance more than content, with a local knowledge base yielding 60 percent follow-through against 10 percent for the same text fetched over the network.
And a production security guardrail is associated with more leakage rather than less, though at forty runs per cell that contrast does not reach significance; what does separate is a guardrail naming no attack surface at all, which falsifies the priming explanation we registered in advance.
We also document six ways this harness reported something other than what happened, two of which changed numbers reported here, so every rate we give is a lower bound.
